
\section{Industry Feedback}

To help improve the super-app ecosystem, we notified our findings super-app developers and the
privacy standards expert teams of the IEEE Standards
Association (IEEE-SA) in China. We summarize their feedback and critical lessons as follows.


\subsection{Notification Process}

\textbf{Methodology.}
Following prior literature, we contacted super-app developers via email addresses extracted from Apple's App Store submissions and super-apps' official websites. In our emails and
online communications, we briefly explained potential leakage risks and
presented our results. We also posed three crucial questions to
the developers and experts: 
(1) whether they were aware of the privacy implications of the data collected under existing legal regulations,
(2) whether they were aware that such information could infer privacy,
and (3) whether they had any remedial plans or suggestions for addressing these privacy risks.


\noindent \textbf{Results.}
We contacted the vendors of all 31 super-apps in Table~\ref{tab:superapp_summary}
and the standards association. We conducted two rounds of
communication. The first round was completed by October 31, 2023.
A second round was initiated for those who did not reply in the first round and was completed by March 31, 2024. Among the 31 vendors, eight acknowledged our report, and four acknowledged and provided feedback on our
questions. 
Most of the unresponsive vendors either didn't reply or
stated that they would forward the report to relevant teams, after which no
further response was received. We speculate that the lack of response from some
vendors could be due to competitive concerns and the protection of business
secrets.
% 
We also engaged the privacy standards department, and they acknowledged our report. Notably, during our communication with company developers, they showed significant interest in our study and frequently asked for details about our attack.


\subsection{Feedback}

Four super-app vendors and the standards association have committed to updating their privacy policies. This update aims to alert users about the recorded mini-app interaction history data during usage and the potential privacy connections. Moreover, we received the following feedback from the developers and experts:

\noindent \textbf{Feedback 1}: The developers acknowledged that ``\textit{We will fix the notification of potential risks of the mini-app interaction history data in our user terms.}'' They further stated that ``\textit{The identified privacy risks indeed exist. Previously, there was a subconscious belief that collecting mini-app interaction history data does not 
violate regulations due to the absence of studies on the privacy leaks they might cause.}''

\noindent \textbf{Feedback 2}: The developers stated: ``\textit{We are working in progress to protect the mini-app interaction history as other sensitive data.}'' Furthermore, they admitted that they had overlooked the privacy issue of mini-app interaction history by saying: ``\textit{We usually follow the latest research papers to protect users' privacy. However, to our knowledge, there is no paper revealing the privacy risks of this data. Thus, we were not aware of the privacy issue.}''

\noindent \textbf{Feedback 3}: The developers promised to ``\textit{improve protection of mini-app interaction history by reducing the time to store such data.}'' However, developers also mentioned the trade-offs between cost and privacy protection: ``\textit{The computation budget to protect privacy is limited. There is a balance between privacy protection and delivering satisfactory services. The user data generated by super-apps is vast, and the cost to protect all data is immeasurable}.''

\noindent \textbf{Feedback 4}: The developers agreed to improve their user term in collecting mini-app interaction history by acknowledging: ``\textit{Before this study, we considered mini-app interaction history as privacy insensitive because we only simply reviewed the data but didn't conduct in-depth research on it.}''

\noindent \textbf{Feedback 5}: The experts from the standards association noted: ``\textit{During our communications with super-app companies, their developers agreed to add mini-app interaction history into the user terms. The association will also standardize the collection and use of such data in the future.}''


\section{Discussions}



\noindent \textbf{Lessons Learned. }
We identified several key lessons to enhance privacy in super-apps. First, developers must rectify all existing misconceptions regarding the boundaries of privacy security. This includes thoroughly re-examining their services to identify and safeguard potentially vulnerable mini-app interaction history. Second, developers must ensure user data privacy rather than merely reacting to privacy breaches. From a regulatory perspective, standards associations should strive to preemptively address new security vulnerabilities instead of waiting to act after breaches occur.

\noindent\revised{\textbf{Potential Interaction Changes.} One potential concern is that the user agreement could introduce bias, influencing volunteers' interaction patterns with the mini-apps. However, we believe this bias has a negligible impact due to the large volume of data, which helps mitigate individual biases. Nevertheless, we acknowledge that measuring bias is challenging since we cannot access ground-truth data (e.g., usage history outside the experimental period or from non-participants).}

\noindent \textbf{Performance of Attacks.}
Note that even though our focus is not on all users, the implications of accurate predictions on a high-confidence subset are profound. This subset, albeit a fraction of the overall users, includes many people due to the colossal scale of super-app users. For instance, for a super-app with over a billion users, accurately inferring privacy attributes of even 1\% users will impact 10 million individuals. This magnitude is substantial and raises serious concerns on personal privacy. Our research highlights this issue and draws attention to robust privacy protection measures, even when only a subset of users is identified.

\noindent \textbf{Privacy Statement Revision.} In the feedback, five super-app vendors promised to revise their user terms. This marks a milestone in super-app privacy, as these industry giants acknowledge and address the privacy implications of mini-app interaction history. By updating the privacy statements, these vendors have demonstrated a commitment to enhancing user privacy and set a precedent for others. The willingness of these vendors to adapt and improve their privacy statements in light of our research is also a strong testament to the generalizability of our findings.

\noindent \textbf{Generalizability.} A key aspect of our research is its broad generalizability, not just to \alipay but to various super-apps as well. The unified interface design and operational paradigms of super-apps and mini-apps lead to the generality of \mini and \op. Therefore, the privacy vulnerabilities we identified from the mini-app interaction history are generalizable findings.




\noindent \textbf{Defense Techniques.} Common defense methods include differential privacy or trusted hardware. However, applying differential privacy to mini-app interaction history could lead to information loss and mislead user interactions with other operations. Utilizing trusted hardware like SGX~\cite{273705} significantly increases costs, which is unsustainable for super-apps with hundreds of millions of active users. Therefore, there is a need to discover more effective defense methods.

